% sec75.tex — §7.5 The reference set: a second, more severe mechanism
% Extracted to avoid encoding corruption in multi-line string substitution

% ----------------------------------------------------------
\subsection{The reference set: a second, more severe mechanism}
\label{sec:refset}
% ----------------------------------------------------------

The results of Section~\ref{sec:lemmaconf} show that Mechanism~1
(implicit re-weighting, $\mathrm{TV}$) orders the three normalisations
within additive aggregation.
A second mechanism governs the fourth cell, and it is more severe.

\paragraph{Mechanism 2: score--evaluator mismatch.}
Proposition~\ref{prop:p1} guarantees Pareto efficiency for the optimal
portfolio only when the \emph{score and the objectives use the same
normalisation}.
When they do not, the guarantee is not weakened but voided.
The observed cost (Table~\ref{tab:fourthcell}: $100\%\to20\%$,
a loss of $80$~pp under optimal assembly) far exceeds what
Mechanism~1 predicts: $\mathrm{TV}=0.107$ implies $\approx9.0$~pp,
accounting for only $14\%$ of the effect.
The gap is the signature of Mechanism~2.

This also explains the interaction of Section~\ref{sec:interaction}.
Under vector normalisation, exact assembly optimises a score that is
misaligned with the evaluator; the $+2.2$~pp gain is residual.
Mechanism~1 (re-weighting) and Mechanism~2 (mismatch) are additive:
mismatch removes the theoretical warrant for optimising, and re-weighting
distorts what is being optimised.

Min--max normalisation satisfies $A_j\equiv1$ over the set on which its
constants are computed.
If that set includes alternatives that cannot be selected, the guarantee
no longer applies: over the $\Nvia$ admissible sites, the attainable
amplitude becomes
\begin{equation}
  \label{eq:refsetA}
  A_j = \frac{\max_{\Nvia}(j)-\min_{\Nvia}(j)}
             {\max_{\Ncand}(j)-\min_{\Ncand}(j)} \;\le\; 1,
\end{equation}
which varies across criteria and re-weights (Mechanism~1 component).
Computing~\eqref{eq:refsetA} for the eleven criteria yields
$\mathrm{TV}=0.107$ (Table~\ref{tab:amprefset}).

\begin{table}[t]
\centering
\caption{Amplitude ratio~\eqref{eq:refsetA} by criterion.}
\label{tab:amprefset}
\begin{tabular}{lcc}
\toprule
Criterion & $A_j$ & $w_j$ \\
\midrule
\texttt{perm}      & 0.482 & 0.152 \\
\texttt{slope}     & 0.683 & 0.113 \\
\texttt{rivers}    & 0.648 & 0.091 \\
\texttt{urban}     & 0.755 & 0.166 \\
\texttt{coast}     & 0.924 & 0.069 \\
\texttt{protected} & 0.999 & 0.066 \\
remaining 5        & 1.000 & ---   \\
\midrule
$\mathrm{TV}$      & \multicolumn{2}{c}{$0.107$} \\
\bottomrule
\end{tabular}
\end{table}

Table~\ref{tab:fourthcell} separates the set over which the
\emph{method} computes its constants from the set over which the
\emph{evaluator} builds the reference set.

\begin{table}[t]
\centering
\caption{Score--evaluator alignment, $Q_2$, SAW min--max, $\kappa=200$.}
\label{tab:fourthcell}
\begin{tabular}{llcc}
\toprule
Method set & Evaluator & greedy & optimal \\
\midrule
$\Ncand$ & $A$  & 68.4\% & \textbf{100.0\%} \\
$\Ncand$ & $A'$ & 15.0\% &  20.0\% \\
$\Nvia$  & $A$  & 46.6\% &  95.0\% \\
$\Nvia$  & $A'$ & 64.6\% & \textbf{100.0\%} \\
\bottomrule
\end{tabular}
\end{table}

Both diagonal cells confirm Proposition~\ref{prop:p1} (optimal assembly
must reach $100\%$ when score and evaluator are aligned).
The off-diagonal costs confirm Mechanism~2: mismatch dominates.
\todo{D2: cite the Cape Verde legal instrument to convert the analytical
  claim into a normative one.}
